\documentclass[11pt]{article}
\usepackage[margin=1in]{geometry}
\usepackage{graphicx}
\usepackage{hyperref}
\usepackage{amsmath}
\usepackage[T1]{fontenc}
\usepackage{amsfonts}
\usepackage{enumitem}
\usepackage{natbib}
\usepackage{listings}
\usepackage{xcolor}
\usepackage{tcolorbox}
\usepackage{titlesec}
\usepackage{fancyhdr}

% Modern color palette inspired by the web design
\definecolor{primaryblue}{RGB}{37,99,235}
\definecolor{lightblue}{RGB}{59,130,246}
\definecolor{darkblue}{RGB}{30,64,175}
\definecolor{accentgray}{RGB}{100,116,139}
\definecolor{lightgray}{RGB}{248,250,252}
\definecolor{mygreen}{rgb}{0,0.6,0}
\definecolor{mygray}{rgb}{0.5,0.5,0.5}
\definecolor{mymauve}{rgb}{0.58,0,0.82}

% Hyperlink setup
\hypersetup{
    colorlinks=true,
    linkcolor=primaryblue,
    urlcolor=lightblue,
    citecolor=darkblue
}

% Section styling
\titleformat{\section}
  {\color{primaryblue}\Large\bfseries}
  {\thesection}{1em}{}
  [\titlerule]

\titleformat{\subsection}
  {\color{lightblue}\large\bfseries}
  {\thesubsection}{1em}{}

% Page styling
\pagestyle{fancy}
\fancyhf{}
\fancyhead[L]{\color{accentgray}\small Graphs in Machine Learning}
\fancyhead[R]{\color{accentgray}\small TD1: Spectral Clustering}
\fancyfoot[C]{\color{accentgray}\thepage}
\renewcommand{\headrulewidth}{0.5pt}
\renewcommand{\headrule}{\hbox to\headwidth{\color{lightblue}\leaders\hrule height \headrulewidth\hfill}}

% Paragraph settings
\setlength\parindent{0pt}
\setlength{\parskip}{0.8\baselineskip}

% Code listing settings
\lstset{
backgroundcolor=\color{lightgray},
basicstyle=\ttfamily\footnotesize,
breakatwhitespace=false,
breaklines=true,
captionpos=b,
commentstyle=\color{mygreen},
escapeinside={\%*}{*)},
extendedchars=true,
frame=single,
frameround=tttt,
rulecolor=\color{lightblue},
keepspaces=true,
keywordstyle=\color{primaryblue}\bfseries,
language=Python,
numbers=left,
numbersep=5pt,
numberstyle=\tiny\color{mygray},
showspaces=false,
showstringspaces=false,
showtabs=false,
stepnumber=1,
stringstyle=\color{mymauve},
tabsize=2
}

% Custom boxes
\newtcolorbox{infobox}{
  colback=lightgray,
  colframe=lightblue,
  fonttitle=\bfseries,
  title=Instructions,
  arc=3mm,
  boxrule=1.5pt
}

\newtcolorbox{questionbox}{
  colback=lightgray,
  colframe=primaryblue,
  fonttitle=\bfseries,
  title=Questions,
  arc=3mm,
  boxrule=1.5pt
}

\begin{document}

% Title page
\begin{center}
    \includegraphics[width=0.4\textwidth]{../../../../../../images/mva/peer_communities2.png}\\[1cm]

    {\Huge\color{primaryblue}\textbf{Practical Session 1}}\\[0.5cm]
    {\LARGE\color{darkblue}Spectral Clustering}\\[1cm]

    {\large\color{accentgray}
    Graphs in Machine Learning\\
    MVA Master Program -- ENS Paris-Saclay}\\[1.5cm]

    {\color{accentgray}
    \textit{michal.valko@inria.fr}\\
    \textit{daniele.calandriello@inria.fr}\\
    \textit{pierre.perrault@inria.fr}\\
    \textit{omar.darwiche-domingues@inria.fr}}
\end{center}

\vfill

\begin{tcolorbox}[colback=lightgray,colframe=lightblue,arc=2mm,boxrule=1pt]
\textbf{About this Session}\\[0.3cm]
In this practical session we will cover fundamental graph building
techniques, and apply them to the Spectral Clustering problem.
The session will be evaluated on a short written report and a final
image segmentation implementation. During
the TD we will implement all the necessary tools to do the
segmentation and answer the
report questions. All the code related to the TD must be submitted,
to provide background for the image segmentation code evaluation.
\end{tcolorbox}

\newpage

\section{Graph Construction}

The file \path{generate_data.py} contains functions that will generate artificial
data for the experiments, as described below. You can run the script to visualize the data.

\begin{infobox}
\textbf{Dataset Types:}
\begin{itemize}[leftmargin=*]
    \item \textbf{$N$-Blob}: Sample random points in $\mathbb{R}^2$ according to $N$
        Gaussian distributions with mean
        $\mu_i = [\cos(2\pi i / N),\sin(2\pi i / N)]$ and variance $\text{Diag}(\sigma^2)$.
    \item \textbf{Two Moons}: Sample random points shaped as two intertwined moons.
    \item \textbf{Point and Circle}: Sample random points from a concentrated Gaussian
        point in the middle and a wide circle around it.
\end{itemize}
\end{infobox}

A prerequisite to build a similarity graph is to define a similarity function
to score the distance between nodes in the graph.
For the rest of the session we will use an inverse exponential function
as the similarity measure, controlled by the Euclidean distance:
\begin{equation}
d(x_i, x_j) = \exp \left\{-\frac{\Vert x_i - x_j \Vert^2_2}{2\sigma^2}\right\}
\end{equation}
The variance $\sigma^2$ controls the bandwidth of the similarity.

\subsection{Implementation Tasks}

In the file \path{build_similarity_graph.py}, complete the following:
\begin{itemize}[leftmargin=*]
	\item  Write the code to build an $\varepsilon$-graph and a
	(OR) $k$-nn graph. Details are in the source code.
	\item Use \path{plot_similarity_graph} to visualize the graph for generated data.
	      The function \path{plot_graph_matrix} in \path{utils.py} may be useful.
	\item Complete the function \path{how_to_choose_epsilon}.
	      You may use \path{min_span_tree} in \path{utils.py}.
\end{itemize}

\begin{questionbox}
\begin{enumerate}[leftmargin=*]
    \item What is the purpose of the option parameter in \path{worst_case_blob}?
    \item  What happens when you change the generating parameter of
        \path{worst_case_blob} in \path{how_to_choose_epsilon}?
        What if the parameter is very large?
    \item Using \path{plot_similarity_graph}, compare $k$-nn to $\varepsilon$ graphs.
          When is it easier to build a connected graph using $k$-nn? When using $\varepsilon$ graphs?
\end{enumerate}
\end{questionbox}

\section{Spectral Clustering}

In the file \path{spectral_clustering.py}, complete:
\begin{itemize}[leftmargin=*]
	\item \path{build_laplacian}
	\item \path{spectral_clustering}
\end{itemize}

\begin{questionbox}
\begin{enumerate}[leftmargin=*]
    \item Build a graph from \path{two_blobs_clustering} data, keeping
        the graph connected. Motivate your eigenvector choices
        and how you computed clustering assignments.
        Compare with built-in $k$-means.
    \item Build a graph from \path{two_blobs_clustering}, but
        make the two components separate.
        How do you choose eigenvectors? Motivate your answer.
\end{enumerate}
\end{questionbox}

\subsection{Adaptive Eigenvector Selection}

Complete \path{spectral_clustering_adaptive}, which uses \path{choose_eig_function}.

\begin{questionbox}
\begin{enumerate}[leftmargin=*,resume]
    \item Look at \path{find_the_bend}. Generate 4-blob data with
        $\sigma^2=0.03$. Plot the first 15 eigenvalues.
        Complete \path{choose_eig_function} to automatically choose
        the number of eigenvectors. The rule must adapt to actual eigenvalues.
    \item Increase variance to $\sigma^2=0.20$ while plotting
        eigenvalues. Use \path{choose_eig_function}.
        Do you see any difference?
    \item When building cluster assignments, did you use thresholding,
        $k$-means or both? When to use each?
    \item What is another use of eigenvalue distributions during clustering,
        beside choosing eigenvectors?
\end{enumerate}
\end{questionbox}

\subsection{Complex Structures}

Complete:
\begin{itemize}[leftmargin=*]
	\item \path{two_moons_clustering}
	\item \path{point_and_circle_clustering}
\end{itemize}

\begin{questionbox}
\begin{enumerate}[leftmargin=*,resume]
    \item  Plot results using spectral clustering and $k$-means in
           \path{two_moons_clustering}. Notice differences?
           Explain considering graph structure.
    \item In \path{point_and_circle_clustering}, compare
        spectral clustering using normal Laplacian $L$ and
        random-walk regularized Laplacian $L_{rw}$.
        Notice differences? Explain considering graph structure.
\end{enumerate}
\end{questionbox}

\subsection{Parameter Sensitivity}

How did you choose parameters when building the graph?
We'll compare clustering solutions while changing parameters.
We use the Adjusted Rand Index \cite{wagner2007comparing} to evaluate assignments,
which outputs values between 0 (unrelated) and 1 (equal).

\begin{questionbox}
\begin{enumerate}[leftmargin=*,resume]
    \item Complete \path{parameter_sensitivity}, and
        generate a plot of the ARI index while varying one parameter
        in graph construction ($\varepsilon$ or $k$).
        Comment on spectral clustering stability.
    \item Without access to \emph{true} labels, how could we evaluate
          clustering results?
\end{enumerate}
\end{questionbox}

\section{Image Segmentation}

Your final task is implementing image segmentation using spectral clustering.
Complete \path{image_segmentation.py}.

\begin{infobox}
\textbf{Pointers:}
\begin{itemize}[leftmargin=*]
    \item Images are $50 \times 50$ pixels RGB in the data folder.
        The image is loaded as $50 \times 50 \times 3$ matrix,
        converted to $2500 \times 3$ where each pixel is an $\mathbb{R}^3$
        vector in RGB space.
    \item Images can be segmented using colors. Build a graph based on
          color distance between pixels and cluster them.
\end{itemize}
\end{infobox}

The function \path{image_segmentation} will plot your segmentation.
Colors might be arbitrary since clustering labels are arbitrary---separation
correctness is what matters.

\begin{questionbox}
\textbf{Final Questions:}
\begin{enumerate}[leftmargin=*]
    \item Document your implementation choices. Well-written, well-commented
        code is sufficient. Include all related code in submission.
    \item A full graph on $50 \times 50$ image has $50^2$ nodes.
        Eigenvalue solving scales as $O(2^{34})$ (seconds to minutes on weak hardware).
        Full HD ($1920 \times 1080$) scales as $O(2^{64})$ (about a month).
        A full graph on millions of nodes requires ~1TB memory.
        Can you think of two simple techniques to reduce computational and memory cost?
    \item Did you use \path{eig} or \path{eigs}? What's the difference?
        How do they scale to large graphs?
\end{enumerate}
\end{questionbox}

\bibliographystyle{plain}
\bibliography{td}

\end{document}
